\chapter*{ABSTRAK}
\addcontentsline{toc}{chapter}{ABSTRAK}

\begin{center}

    \textbf{Modifikasi \textit{Framework} TOGAF untuk Kementerian Pertahanan RI\\yang Berbasis pada C6ISR} \\
    \vspace{0.4cm}
    oleh \\
    \vspace{0.4cm}
    Marcel Suandi Tambing\\
    23524316\\
\end{center}


\begin{spacing}{1.0}
Kementerian Pertahanan Republik Indonesia (Kemhan RI) wajib menyusun Arsitektur Sistem Pemerintahan Berbasis Elektronik (SPBE) sebagai kementerian, sekaligus perlu menjangkau kapabilitas operasional sebagai institusi pertahanan. Namun, sistem di lingkungan Kemhan RI masih berkembang secara terpisah (\textit{silo}) pada lapisan bisnis dan tata kelola, data dan aplikasi, teknologi, serta operasional. Framework \textit{enterprise architecture} (EA) pertahanan, seperti DoDAF, MODAF, dan NAF, kuat dalam muatan pertahanan, tetapi berfokus pada deskripsi arsitektur atau mengikuti kebijakan NATO, sedangkan TOGAF unggul dalam metode, tetapi tidak menyediakan muatan pertahanan. Akibatnya, belum ada framework EA yang memenuhi kedua kebutuhan tersebut.

Penelitian ini bertujuan memodifikasi TOGAF \textit{Architecture Development Method} (ADM) berbasis C6ISR (\textit{Command, Control, Communications, Computers, Cyber-Defense, Combat Systems, Intelligence, Surveillance, and Reconnaissance}), menerapkannya pada Kemhan RI, serta mengevaluasi kelayakan dan kontribusinya. Penelitian menggunakan \textit{Design Science Research Methodology} (DSRM). Dari sepuluh kebutuhan framework yang diturunkan dari domain pertahanan, regulasi nasional, dan praktik terbaik, empat kebutuhan belum dipenuhi TOGAF ADM standar, yaitu akomodasi domain C6ISR, keamanan siber sejak perancangan, interoperabilitas lintas entitas, serta keterkaitan dengan doktrin dan strategi pertahanan.

Hasil penelitian berupa TOGAF Pertahanan Indonesia, yaitu TOGAF ADM yang mempertahankan siklus pengembangannya, menambahkan Fase C6ISR \textit{Architecture} sebagai domain arsitektur tersendiri, menyesuaikan muatan lima fase lainnya, serta menetapkan sembilan prinsip arsitektur. Penerapannya pada Kemhan RI dari Fase \textit{Preliminary} hingga \textit{Migration Planning} menghasilkan arsitektur transformasi digital pada lima domain, yaitu bisnis, data, aplikasi, teknologi, dan C6ISR. Arsitektur tersebut antara lain mengusulkan unit \textit{Digital Transformation Officer}, Portal Satu Data Pertahanan, rasionalisasi 71 aplikasi khusus menjadi 41 aplikasi terintegrasi, \textit{Data Center} Utama dan \textit{Disaster Recovery Center} aktif-aktif, serta Tampilan Pemantauan Terpusat lintas matra. Kesenjangan arsitektur dihimpun menjadi delapan paket pekerjaan dalam \textit{roadmap} 2026--2030 dengan tiga tahap migrasi dan delapan indikator dampak.

\vspace{0.4cm}
\noindent\textbf{Kata kunci:} Arsitektur \textit{enterprise}, C6ISR, Kementerian Pertahanan, TOGAF, transformasi digital.
\end{spacing}
